\chapter*{Kế hoạch thực hiện công việc nhóm}
\addcontentsline{toc}{chapter}{Kế hoạch thực hiện công việc nhóm}
\markboth{KẾ HOẠCH THỰC HIỆN CÔNG VIỆC NHÓM}{}
\setlength{\parskip}{6pt}

Bảng dưới đây tóm tắt kế hoạch, phân công và mức hoàn thành công việc theo các hạng mục chuyên môn của đề tài. Nhóm gồm hai thành viên: \textbf{SV1} --- \studentname{}; \textbf{SV2} --- \studentnametwo{}.

\begin{table}[H]
\centering
\caption{Kế hoạch và phân công thực hiện công việc}
\label{tab:ke-hoach-nhom}
\begin{tabularx}{\linewidth}{@{}p{0.04\linewidth}X p{0.16\linewidth}p{0.15\linewidth}@{}}
\toprule
\textbf{TT} & \textbf{Nội dung} & \textbf{Phụ trách} & \textbf{Mức độ} \\
\midrule
1 & Khảo sát bài toán, nghiên cứu nền tảng an toàn LLM/RAG & Cả hai & Hoàn thành \\
2 & Phân tích yêu cầu FR/NFR, mô hình mối đe doạ, kiến trúc & Cả hai & Hoàn thành \\
3 & Hiện thực gateway và các lớp guard (Input, RAG Context, Output/DLP) & SV1 & Hoàn thành \\
4 & Truy hồi lai tối ưu (BM25 + embedding, RBAC, vectorized) & SV1 & Hoàn thành \\
5 & Kho tri thức doanh nghiệp tổng hợp và khung so sánh A/B & SV1 & Hoàn thành \\
6 & Bộ đánh giá bảo mật tổng hợp và chạy đo chỉ số & SV2 & Hoàn thành \\
7 & Tối ưu lớp luật chống mạo danh và phân tích hiệu quả tường & SV2 & Hoàn thành \\
8 & Tích hợp Semantic Judge và thí nghiệm đối chiếu mô hình cục bộ & SV2 & Hoàn thành \\
9 & Kiểm thử tự động, chỉnh sửa theo kết quả & Cả hai & Hoàn thành \\
10 & Soạn thảo và hoàn thiện báo cáo & Cả hai & Hoàn thành \\
\bottomrule
\end{tabularx}
\end{table}

\textbf{Nguyên tắc làm việc:} hai thành viên rà soát chéo kết quả của nhau trước khi hợp nhất; mọi số liệu đánh giá bám sát artifact/script tái lập được; không mở rộng kết luận vượt quá bằng chứng thực nghiệm.
